% 存放页边距、行距、页眉等格式设置
%\xeCJKsetup{EmboldenFactor=1}
%\setCJKfamilyfont{zhsong}{SimSun}[AutoFakeBold]

\geometry{ 
	a4paper, 
	left=3cm, right=2.5cm, 
	top=2cm, bottom=2cm, 
}

% 设置粗宋体和粗黑体

\newCJKfontfamily{\cuheiti}[AutoFakeBold={3.17}]{SimHei}

\newcommand{\yihao}{\fontsize{26pt}{26pt}\selectfont}       % 一号, 单倍行距
\newcommand{\xiaoyi}{\fontsize{24pt}{24pt}\selectfont}      % 小一, 单倍行距
\newcommand{\erhao}{\fontsize{22pt}{1.25\baselineskip}\selectfont}       % 二号, 1.25倍行距
\newcommand{\xiaoer}{\fontsize{18pt}{18pt}\selectfont}      % 小二, 单倍行距
\newcommand{\sanhao}{\fontsize{16pt}{16pt}\selectfont}      % 三号, 单倍行距
\newcommand{\xiaosan}{\fontsize{15pt}{15pt}\selectfont}     % 小三, 单倍行距
\newcommand{\sihao}{\fontsize{14pt}{16.1pt}\selectfont}       % 四号, 单倍行距
\newcommand{\xiaosi}{\fontsize{12pt}{13.8pt}\selectfont}      % 小四, 1.15倍行距
\newcommand{\xiaoxiaosi}{\fontsize{11pt}{12.65pt}\selectfont}      % 小四, 1.15倍行距
\newcommand{\wuhao}{\fontsize{10.5pt}{10.5pt}\selectfont}   % 五号, 单倍行距
\newcommand{\xiaowu}{\fontsize{9pt}{9pt}\selectfont}        % 小五, 单倍行距

\newcommand{\upcite}[1]{\textsuperscript{\textsuperscript{\cite{#1}}}} % 上角标引用参考文献

%配置ctex样式
\ctexset {
	part/pagestyle = empty,
	chapter/pagestyle = fancy,
	chapter = {
		name={,.},
		number=\textbf{\arabic{chapter}},
		format=\xiaosi\fontspec{Arial}\bfseries\raggedright, % 追加跳转
		beforeskip = {0pt},
		afterskip={10pt},
		fixskip=true,
	},
	section = {
		name={,},
		% number=\textbf{\thechapter.\arabic{section}},
		number=\arabic{chapter}.\arabic{section},
		format=\xiaoxiaosi\fontspec{Arial}\bfseries\raggedright, % 追加跳转
		beforeskip = {0pt},
		afterskip={10pt},
		fixskip=true,
	},
	subsection = {
		name={,},
		number=\textbf{\thesection.\arabic{subsection}},
		format=\xiaoxiaosi\fontspec{Arial}\bfseries\raggedright, % 追加跳转
		beforeskip = {0pt},
		afterskip={10pt},
		fixskip=true,
	},
	subsubsection = {
		number=\textbf{\alpha{subsubsection}},
		format=\xiaoxiaosi\fontspec{Arial}\bfseries\raggedright,
		beforeskip = {0pt},
		afterskip={10pt},
		fixskip=true,
	},
  	listfigurename=List of Figure,
	listtablename=List of Table,
	figurename=Figure,
	tablename=Table,
	bibname=Reference,
	indexname=Index,
}
% 摘要页的页眉设置
\fancypagestyle{toastyle}
{
	\fancyhf{} % 清空
	\fancyhead[C]{
  		\ifodd\value{page}
     		\SansC\xiaowu Abstract
  		\else
     		\SansC\xiaowu Xi'an University of Posts and Telecommunications Final Year Project (Paper)
  		\fi
	}
	\fancyfoot[C]{\thepage}
}
% 目录页的页眉设置
\fancypagestyle{tocstyle}
{
    \fancyhf{} % 清空
    \fancyhead[C]{
  		\ifodd\value{page}
     		\SansC\xiaowu Contents
  		\else
     		\SansC\xiaowu Xi'an University of Posts and Telecommunications Final Year Project (Paper)
  		\fi
}
    \fancyfoot[C]{\thepage}
}
%图目录页的页眉设置
\fancypagestyle{lofstyle}
{
	\fancyhf{} % 清空
	\fancyhead[C]{
  		\ifodd\value{page}
	 		\SansC\xiaowu List of Figure
  		\else
	 		\SansC\xiaowu Xi'an University of Posts and Telecommunications Final Year Project (Paper)
  		\fi
	}
	\fancyfoot[C]{\thepage}
}
%表目录页的页眉设置
\fancypagestyle{lotstyle}
{
	\fancyhf{} % 清空
	\fancyhead[C]{
  		\ifodd\value{page}
	 		\SansC\xiaowu List of Table
  		\else
	 		\SansC\xiaowu Xi'an University of Posts and Telecommunications Final Year Project (Paper)
  		\fi
	}
	\fancyfoot[C]{\thepage}
}
% 正文的页眉设置
% —— 可回退的英文字体宏（优先用系统字体，找不到就退而求其次）——
\IfFontExistsTF{Arial}{%
  \newcommand{\SansA}{\fontspec{Arial}}%
}{%
  \IfFontExistsTF{TeX Gyre Heros}{\newcommand{\SansA}{\fontspec{TeX Gyre Heros}}}{\newcommand{\SansA}{\sffamily}}%
}
\IfFontExistsTF{Calibri}{%
  \newcommand{\SansC}{\fontspec{Calibri}}%
}{%
  \IfFontExistsTF{Carlito}{\newcommand{\SansC}{\fontspec{Carlito}}}{\newcommand{\SansC}{\SansA}}%
}
\renewcommand{\chaptermark}[1]{%
  \markboth{\chaptername~\thechapter.\ #1}{}%
}
% —— 页眉：偶数页固定标题，奇数页章节标题（Calibri 不在则回退）——
\fancypagestyle{fancy}{
    \fancyhead[R]{}                % 清空右页眉
	\fancyhead[L]{}                % 清空左页眉
	\fancyfoot[l]{} 			   % 清空左页脚
	\fancyfoot[r]{} 			   % 清空右页脚
    \fancyhead[C]{
  	\ifodd\value{page}
     	{\SansC\xiaowu\nouppercase{\leftmark}}
  	\else
     	\SansC\xiaowu Xi'an University of Posts and Telecommunications Final Year Project (Paper)
  	\fi
}
    \fancyfoot[C]{\footnotesize\thepage}
    \renewcommand{\headrulewidth}{0.75pt}
}
% ===== 公式编号格式设置 =====
\renewcommand{\theequation}{\thechapter-\arabic{equation}}

% ===== 图片编号格式设置 =====
\renewcommand{\thefigure}{\thechapter.\arabic{figure}}

% ===== 表格编号格式设置 =====
\renewcommand{\thetable}{\thechapter.\arabic{table}}

% ===== 图表标题格式设置 =====
\captionsetup[figure]{
    font={normalsize},           % 标题字号与正文相同（11pt）
    labelfont={},              % "Fig 2-1" 加粗
    textfont={},                 % 标题文字常规
    labelsep=space,              % 编号与标题用空格分隔
    skip=10pt,                   % 标题与图片间距10pt
    position=below,              % 标题在图片下方
    name={Figure}                   % 使用 "Figure" 而非 "图"
}
\captionsetup[table]{
    font={normalsize},
    labelfont={},
    textfont={},
    labelsep=space,
    skip=10pt,
    position=above,              % 表格标题在上方
    name={Table}
}